% Part: lambda-calculus
% Chapter: church-rosser
% Section: beta-eta-reduction

\documentclass[../../../include/open-logic-section]{subfiles}

\begin{document}

\olfileid{lam}{cr}{be}

\olsection{$\beta\eta$-తగ్గింపు}

$\beta\eta$-తగ్గింపుకు ($\bered$) చర్చ్--రోసర్ లక్షణం ఉంది.

\begin{lem}\ollabel{lem:one-par}
  $M \beredone M'$ అయితే $M \beredpar M'$.
  \sourcecorrection{OLTELAMCRBE-001}{మూలంలో $\beredone$ సంకేతాన్ని ఈ ఉపసిద్ధాంతంలో వాడారు; ముందు $\beta\eta$ విభాగం మాత్రం $\bredone$, $\eredone$లను కలిగి ఉండే స్వప్రావర్తక-సంక్రమణ సంబంధం $\bered$ను మాత్రమే నిర్వచించింది. $\beredone$కు టెక్స్ మాక్రో ఉన్నా, అది ఏ ఒక-దశ సంబంధమో ప్రత్యేక గణిత నిర్వచనం అక్కడ లేదు. మూల సూత్రాన్ని అలాగే ఉంచి ఈ అర్థ-నిర్వచన లోటును ప్రకటిస్తున్నాం.}
\end{lem}

\begin{proof}
  $M \beredone M'$ యొక్క \tetoken{వ్యుత్పత్తి}{!!{derivation}}పై
  ఆగమన పద్ధతిలో నిరూపిస్తామని మూలం చెబుతుంది.
  $M \eredone M'$ అనేది $\eta$-పరివర్తనం ద్వారా అయితే
  (అంటే \olref[syn][eta]{defn:beredone}),
  \olref[pbe]{thm:refl}ను వాడుతుంది. మిగతా సందర్భాలు
  \olref[b]{lem:one-par}లోని వాదనల మాదిరిగానే అని చెబుతుంది.
  \sourcecorrection{OLTELAMCRBE-002}{మూలంలో “$M \bredone M'$ అనేది $\eta$-పరివర్తనం ద్వారా” అని ఉంది; కానీ $\bredone$ బీటా సంకోచనం, ఏటా సంకోచనానికి ముందరి నిర్వచనం $\eredone$. తెలుగు వాక్యంలో ఆ ఒక్క బాణాన్ని సరిచేశాం. మూలం పేర్కొన్న \olref[pbe]{thm:refl} ఒక్కటే సందర్భంలో ఏటా సంకోచనానికి సరిపోదు; అయిదవ సమాంతర నియమమూ కావాలి. అలాగే “మిగతా సందర్భాలు” అన్న ఆధార \olref[b]{lem:one-par}లో OLTELAMCRB-001 అనుకూల-సందర్భ ఖాళీ ఇంకా ఉంది. ఇక్కడ పూర్తి కొత్త నిరూపణను ప్రకటించడం లేదు.}
\end{proof}

\begin{lem}\ollabel{lem:par-red}
  $M \beredpar M'$ అయితే $M \bered M'$.
\end{lem}

\begin{proof} $M \beredpar M'$ యొక్క
  \tetoken{వ్యుత్పత్తి}{!!{derivation}}పై ఆగమన పద్ధతిలో నిరూపిస్తాం.

  చివరి నియమం \olref[pbe]{defn:beredpar5} అయితే,
  $M$ అనేది $\lambd[x][Nx]$, $M'$ అనేది $N'$;
  ఇక్కడ $x$, $N$, $N'$ తగిన చరం, పదాలు;
  $x \notin FV(N)$, $N \beredpar N'$. కనుక ముందుగా
  $\eta$-పరివర్తనంతో $\lambd[x][Nx]$ను $N$కు
  తగ్గించి, ఆ తరువాత $N \bered N'$ అని చూపే
  $\beredone$ దశల క్రమాన్ని అనుసరించవచ్చు;
  రెండవది ఆగమన పరికల్పన వల్ల వస్తుంది.
  \sourcecorrection{OLTELAMCRBE-003}{మూలం చివరి అయిదవ నియమపు సందర్భాన్నే వివరించి, తొలి నాలుగు సందర్భాలను రాయలేదు. అందులో వాడిన $\beredone$ ఒక-దశ సంకేతానికి OLTELAMCRBE-001లో తెలిపినట్లు ప్రత్యేక నిర్వచనం లేదు; అందువల్ల “$N \bered N'$ను చూపే $\beredone$ దశలు” అని మూలం చేసిన గుర్తింపును పూర్తి స్థాపనగా తీసుకోలేదు. ఈ నిరూపణ భాగాన్ని ముద్రిత పరిమితితోనే ఉంచాం.}
\end{proof}

\begin{lem}\ollabel{lem:str}
  $\beredpar$ను కలిగి ఉండే కనిష్ఠ సంక్రమణ సంబంధం $\bered$.
\end{lem}

\begin{proof}
  \olref[b]{lem:str}లో మాదిరిగానే.
\end{proof}

\begin{thm}\ollabel{thm:cr}
  $\bered$ చర్చ్--రోసర్ లక్షణాన్ని తీరుస్తుంది.
\end{thm}

\begin{proof}
  \olref[dap]{thm:str}, \olref[pbe]{thm:cr},
  \olref{lem:str} ప్రకారం వస్తుంది.
  \sourcecorrection{OLTELAMCRBE-004}{ఈ తుది వాదన ముందరి సమాంతర $\beta\eta$ చర్చ్--రోసర్ సిద్ధాంతం, ఇక్కడి కనిష్ఠ సంక్రమణ ఉపసిద్ధాంతంపై ఆధారపడుతుంది. OLTELAMCRPBE-002–005, OLTELAMCRBE-001–003గా ప్రకటించిన నిర్వచన, నిరూపణ పరిమితులు ఇంకా పూరించలేదు. మూల సూచనలను నిలిపాం; దీన్ని స్వతంత్రంగా ధ్రువీకరించిన పూర్తి నిరూపణగా పేర్కొనడం లేదు.}
\end{proof}
\end{document}
